% !TEX root = ../main.tex
% ---------------------------------------------------------------------
\section{Conception}
\label{sec-conception}
% ---------------------------------------------------------------------

Chaque étape est justifiée par les propriétés de la \cref{sec-theorie} ; les valeurs finales
des paramètres sont rassemblées au \cref{tab-parametres}.

\subsection{Carte de constellation (Q3)}
\label{sec-pics}

\paragraph{Définition.} Pour un voisinage $V(m,k)$ de $(2\mv+1)$ trames sur $(2\kv+1)$
canaux centré en $(m,k)$, la carte de constellation est
\begin{equation}
  \mathcal{P} = \Bigl\{ (m,k) : S[m,k] = \max_{V(m,k)} S \ \text{ et } \ \SdB[m,k] \ge \tau[m,k] \Bigr\},
  \tag{4}
\end{equation}
où $\tau[m,k]$ est un seuil \emph{local}. D'après la \cref{sec-maxima}, cette définition est
exactement invariante au gain, et invariante au filtrage sous la condition (5).

\paragraph{Voisinage de départ.} Avec $\mv = 3$ et $\kv = 5$, un pic domine
$\pm 3 \times \qty{64}{\milli\second} = \pm\qty{192}{\milli\second}$ et
$\pm 5 \times \qty{7.8}{\hertz} = \pm\qty{39}{\hertz}$. La demi-largeur en fréquence vaut
un demi-ton vers \qty{650}{\hertz} : elle est assez étroite pour que $\Delta g(V)$
reste petit, et assez large pour qu'un même lobe spectral (largeur
$\Delta f \approx \qty{11}{\hertz}$, soit 1,4 canal) ne donne qu'un seul pic.

\paragraph{Seuil local et quota par bande.} Le seuil $\tau$ est la médiane de $\SdB$ dans la
bande de fréquence, augmentée d'une marge : il suit le gain et le niveau moyen de la bande,
donc il est invariant au gain et quasi invariant à un filtrage lent. On impose en plus un
nombre fixe de pics par bande et par seconde, en gardant les plus forts : la densité $\rho$
est ainsi maîtrisée, quels que soient le volume et le genre. Les 6 bandes sont en
progression géométrique, $f_b = f_{\min}(f_{\max}/f_{\min})^{b/N_b}$ entre 200 et
\qty{4000}{\hertz}, de rapport constant $20^{1/6} \approx \num{1.65}$ (environ $\num{0.72}$
octave, soit 8,6 demi-tons par bande) : chaque bande couvre le même intervalle musical. Un
découpage linéaire placerait la moitié des pics imposés par le quota au-dessus de
\qty{2.1}{\kilo\hertz}, là où l'énergie musicale est faible et les pics fragiles.

\paragraph{Question 3 : pourquoi $\rho p^2$ et pas $p$ seul.} Notons $\rho$ la densité de
pics (pics/s) et $p$ la probabilité qu'un pic survive à la dégradation. Une requête de
durée $\Tq$ contient environ $\rho\Tq$ ancres ; chacune forme $F$ paires, et une paire
n'est retrouvée que si ses deux pics survivent. Si les survies des deux pics sont
indépendantes, une clé survit avec la probabilité $p^2$ et le nombre de clés correctes
retrouvées vaut
\begin{equation}
  \E[\text{clés correctes}] = F\,\rho\,\Tq\,p^2 .
  \label{eq-cles-correctes}
\end{equation}
À $F$ et $\Tq$ fixés, c'est donc $\rho p^2$ qui mesure ce qu'un réglage du voisinage apporte
au score du bon morceau. Le critère $p$ seul est trompeur, car $p$ et $\rho$ varient en sens
inverse : un voisinage \emph{large} ne garde que des pics très dominants, qui survivent bien
($p$ grand) mais sont rares ($\rho$ petit) ; un voisinage \emph{étroit} donne beaucoup de pics,
dont une bonne partie, faiblement dominants, sont déplacés par le bruit ($p$ petit).
Maximiser $p$ conduirait à une empreinte presque vide. Le carré compte aussi : une clé exige
la survie de \emph{deux} pics, donc une baisse de $p$ pèse plus qu'une hausse équivalente de
$\rho$. L'hypothèse d'indépendance est d'ailleurs plutôt pessimiste : deux pics proches
subissent le même bruit local, et une clé survit alors avec une probabilité supérieure à
$p^2$.

\paragraph{Mesure.} Avec le voisinage de départ ($\mv = 3$, $\kv = 5$), le quota de
30~pics/s et une marge de \qty{6}{\decibel}, on obtient $\rho \approx 25$~pics/s : le quota
n'est pas atteint dans les passages où peu de maxima dépassent le seuil local. Nous avons balayé le voisinage sans quota, de $\pm 1$ à $\pm 5$ trames
et de $\pm 1$ à $\pm 12$ canaux, sur 100 extraits de \qty{8}{\second} (un par morceau)
bruités à \qty{10}{\decibel} (\cref{fig-p-rho}). Comme prévu, $\rho$ et $p$ varient en sens
inverse : quand le voisinage grandit, $\rho$ chute de 308 à 23~pics/s, et $p$ ne monte que
de \num{0.30} à \num{0.42}. $\rho p^2$ décroît donc de \num{27.7} à \num{4.1}. Le réglage
retenu ($\pm 3$ trames $\times$ $\pm 5$ canaux, avec quota) donne $\rho = 26$~pics/s,
$p = \num{0.55}$ et $\rho p^2 = \num{7.8}$, soit près de deux fois plus que le voisinage
sans quota de même densité ($p = \num{0.42}$, $\rho p^2 = \num{4.1}$) : le quota garde les
pics les plus forts de chaque bande, qui sont ceux qui dépassent le plus le bruit.
Maximiser $\rho p^2$ pousserait vers des voisinages minuscules, mais la taille de l'index
croît comme $\rho F$ et les appariements fortuits comme $\rho^2$ : on fixe donc $\rho$ dans
la plage du sujet (25 à 35~pics/s) par le quota, et c'est à densité fixée que le critère
$\rho p^2$ départage les réglages.

\begin{figure}[t]
  \centering
  \inclurefig[0.8\linewidth]{p_en_fonction_de_rho}{p sous bruit à 10 dB en fonction de rho, un point par voisinage}
  \caption{Proportion $p$ de pics conservés après ajout d'un bruit blanc à
  \qty{10}{\decibel}. Voisinages de $\pm 1$, $\pm 3$ et $\pm 5$ trames sur 3 à 25 canaux,
  sans quota ; étoile : réglage retenu, avec quota. 100 extraits de \qty{8}{\second}.
  (a)~$\rho$ en fonction de la largeur du voisinage ; (b)~$p$ en fonction de $\rho$ ;
  (c)~critère $\rho p^2$. $p$ augmente quand le voisinage grandit, mais beaucoup moins vite
  que $\rho$ ne diminue ; à densité égale, le quota double $\rho p^2$.}
  \label{fig-p-rho}
\end{figure}

\subsection{Empreintes par paires de pics (Q4)}
\label{sec-empreintes}

\paragraph{La clé.} À chaque pic ancre $(m_1,k_1)$ on associe les $F$ premiers pics cibles
$(m_2,k_2)$ de la zone $\dm_{\min} \le m_2 - m_1 \le \dm_{\max}$, $|k_2 - k_1| \le \dk_{\max}$,
et chaque paire donne la clé
\begin{equation}
  u = \bigl(Q_f(k_1),\, Q_f(k_2),\, \dm\bigr), \qquad \dm = m_2 - m_1, \qquad Q_f(k) = \lfloor k/\gf \rfloor ,
  \tag{6}
\end{equation}
rangée dans l'index avec la valeur (morceau, $m_1$). La clé ne dépend que de positions de
pics, donc elle est invariante au gain ; elle ne contient que l'écart $m_2 - m_1$, qu'une
translation (propriété (1)) laisse inchangé, et c'est pour cela que l'instant absolu $m_1$
est rangé dans la valeur et non dans la clé. Le triplet est codé en un seul entier
$u = (Q_f(k_1)\,\Nf + Q_f(k_2))\,\Nt + (\dm - \dm_{\min})$, inférieur à
$M = \Nf^2\Nt \le 2^{32}$.

\paragraph{Zone cible.} À \num{15.6} trames/s, $\Delta t \in [\num{0.1} ; 2]$\,s donne
$\dm \in [2, 31]$, soit $\Nt = 30$ valeurs. On exclut $\dm = 0$ (pics d'un même accord, très
fréquents, peu spécifiques) et $\dm = 1$ (où la migration d'un pic d'une trame change le plus
souvent la clé). La borne $\dk_{\max} = 60$ canaux (\qty{470}{\hertz}) garde la cible dans la
même région musicale. La zone contient en moyenne
$30 \times \num{1.9} \times 121/490 \approx 14$ pics ; $F$ doit rester en dessous pour que les
ancres des bords de bande aient assez de cibles.

\paragraph{Pourquoi des paires.} Une clé « 1 pic » survit avec la probabilité $p$, une clé
« 1 paire » avec $p^2$, mais il y a $F$ fois plus de paires que de pics : le nombre de clés
correctes est multiplié par $Fp$, et les paires gagnent dès que $F\,p > 1$. Dans le même
temps, le nombre de clés possibles passe de $K = 513$ à $M = \Nf^2 \Nt \approx \num{7.9e6}$ :
les paires sont plus nombreuses à survivre \emph{et} \num{15000} fois plus discriminantes, au
prix d'un index $F$ fois plus gros (\cref{tab-pic-paire}).

\begin{table}[h]
  \centering\small
  \caption{Clés « 1 pic » contre clés « 1 paire », mesurées sur 100 requêtes de
  $\Tq = \qty{8}{\second}$ alignées, $\rho = \num{26.1}$~pics/s, $F = 10$. Colonne de gauche :
  exemple de l'énoncé ($\Tq = \qty{8}{\second}$, $\rho = 30$, $F = 10$, $p = \num{0.35}$).
  Nombre de paires mesuré : \num{1707}, et non $F\rho\Tq = \num{2090}$, car les ancres de fin
  d'extrait n'ont pas toutes leurs $F$ cibles.}
  \label{tab-pic-paire}
  \begin{tabular}{@{}lccccc@{}}
    \toprule
    & & \multicolumn{2}{c}{Bruit blanc \qty{10}{\decibel}} & \multicolumn{2}{c}{MP3 \qty{64}{\kilo\bit\per\second}} \\
    \cmidrule(lr){3-4}\cmidrule(l){5-6}
    & Énoncé (pic / paire) & 1 pic & 1 paire & 1 pic & 1 paire \\
    \midrule
    Clés dans la requête
      & 240 / 2400 & 209 & \num{1707} & 209 & \num{1707} \\
    Probabilité de survie : $p$ / $p^2$
      & \num{0.35} / \num{0.12} & \num{0.56} & \num{0.31} & \num{0.86} & \num{0.73} \\
    Survie des clés mesurée
      & -- & -- & \num{0.36} & -- & \num{0.75} \\
    Clés correctes retrouvées (mesurées)
      & 84 / 294 & 116 & 612 & 179 & \num{1276} \\
    Clés possibles : $K$ / $M$
      & 513 / \num{7.9e6} & 513 & \num{7.9e6} & 513 & \num{7.9e6} \\
    \bottomrule
  \end{tabular}
\end{table}

La condition $Fp > 1$ est largement remplie : $Fp = \num{5.6}$ sous bruit à
\qty{10}{\decibel} et \num{8.6} en MP3. Les paires donnent 4 à 7 fois plus de clés correctes
que les pics isolés (612 contre 116, \num{1276} contre 179), tout en étant \num{15000} fois plus
discriminantes. La survie des clés mesurée dépasse un peu $p^2$ (\num{0.36} contre
\num{0.31} sous bruit) : les deux pics d'une paire ne disparaissent pas indépendamment, les
pics forts survivent ensemble et les pics faibles disparaissent ensemble. L'hypothèse
d'indépendance est donc légèrement pessimiste, comme annoncé à la \cref{sec-pics}.

\paragraph{Regroupement en fréquence : $\gf$.} Regrouper $\gf$ canaux voisins en une classe
donne $\Nf = \lceil K/\gf \rceil$ classes et $M = \Nf^2\Nt$ clés possibles. Si les clés étaient
réparties uniformément, une requête de $\nq$ clés interrogée contre un index de $\Nu$ entrées
produirait en moyenne $\nq\Nu/M$ appariements fortuits (\cref{sec-decision}). Le
\cref{tab-gammaf} donne ces grandeurs pour $\gf = 1, 2, 4, 8$.

\begin{table}[h]
  \centering\small
  \caption{Effet du regroupement $\gf$. Parasites attendus : modèle uniforme $\nq\Nu/M$, avec
  $\nq \approx \num{1015}$ clés mesurées par requête de \qty{5}{\second} et
  $\Nu = \num{677171}$ (notre index de 100 extraits de \qty{30}{\second}). Plancher : seuil
  prévu par le modèle uniforme (\cref{eq-plancher}). $\sthr$ mesuré : seuil recalibré pour
  chaque $\gf$ sur 50 négatives (plage selon le RSB). \TBD{} : $\Tq = \qty{5}{\second}$,
  200 requêtes positives ; $\TFD = 0/50$ partout.}
  \label{tab-gammaf}
  \begin{tabular}{@{}rrrrrrrr@{}}
    \toprule
    $\gf$ & $\Nf$ & $M$ & Parasites & Plancher & $\sthr$ mesuré & \TBD{} \qty{10}{\decibel} & \TBD{} \qty{0}{\decibel} \\
    \midrule
    1 & 513 & \num{7.9e6} & 87         & 3 & 21 à 26 & \qty{95.5}{\percent} & \qty{81.5}{\percent} \\
    2 & 257 & \num{2.0e6} & 347        & 3 & 26 à 33 & \qty{96.0}{\percent} & \qty{81.5}{\percent} \\
    4 & 129 & \num{5.0e5} & \num{1380} & 4 & 35 à 41 & \qty{96.0}{\percent} & \qty{80.5}{\percent} \\
    8 & 65  & \num{1.3e5} & \num{5430} & 6 & 59      & \qty{94.0}{\percent} & \qty{71.0}{\percent} \\
    \bottomrule
  \end{tabular}
\end{table}

\paragraph{Question 4 : le meilleur $\gf$ n'est pas celui qui donne le moins de parasites.}
Les parasites croissent comme $\gf^2$, donc $\gf = 1$ les minimise. Mais deux effets jouent
en sens inverse. D'abord, les parasites gênent peu tant qu'ils restent dispersés : la
décision porte sur la hauteur du plus grand pic de l'histogramme des $\delta$, pas sur le
nombre total d'appariements, et le maximum de nombreuses variables de Poisson de petite
moyenne croît très lentement (le plancher de seuil ne passe que de 3 à 6 quand les
parasites sont multipliés par 64). Ensuite, un regroupement absorbe les petits déplacements
de pics : si un pic survivant se décale d'un canal avec la probabilité $j$, il ne change de
classe qu'avec la probabilité $j/\gf$, et une clé survit avec la probabilité
$\bigl[p\,(1 - j/\gf)\bigr]^2$. Pour $j = \num{0.3}$, passer de $\gf = 1$ à 2 multiplie les clés
correctes par environ \num{1.5}. Le bon critère est donc la \emph{séparation} entre le score
du bon morceau et celui des mauvais. Quand $\gf$ devient grand, une même clé se répète dans
chaque morceau, les votes fortuits s'empilent dans les mêmes cases et la spécificité
s'effondre : l'optimum est intermédiaire.

\paragraph{Mesure.} Le \cref{tab-gammaf} confirme les deux effets, mais le second est bien
plus fort que prévu : le score médian du bon morceau à \qty{0}{\decibel} augmente avec $\gf$
(49, 65, 76, 88), mais le seuil nécessaire pour garder $\TFD = 0$ monte de 21 à 59, et non de
3 à 6, car les clés réelles sont loin d'être uniformes (\cref{sec-decision}). D'où un
plateau : $\gf = 1$ et 2 donnent le même \TBD{}, puis le \TBD{} à \qty{0}{\decibel} chute à
\qty{71}{\percent} pour $\gf = 8$ (à \qty{10}{\decibel}, les écarts ne valent que 1 à 4
requêtes sur 200). Nous gardons $\gf = 1$ : même \TBD{} que $\gf = 2$, quatre fois moins
d'appariements fortuits et le seuil le plus bas.

\paragraph{Choix de $F$.} Avec les mêmes requêtes, $F = 5$, 10 et 15 donnent
\qtylist{78.0;81.5;81.5}{\percent} de \TBD{} à \qty{0}{\decibel}, pour un index de
\numlist{368190;677171;905862} entrées. Le score augmente avec $F$, mais le seuil aussi :
le \TBD{} plafonne à partir de $F = 10$, alors que l'index continue de croître. $F = 10$ est
le plus petit $F$ qui atteint ce plafond.

\subsection{Index inversé}
\label{sec-index}

L'index inversé range, pour chaque clé, la liste des couples (morceau, $m_1$) : une requête
ne lit que ses $\nq$ clés au lieu de se comparer à chaque morceau. Nous le stockons comme un
tableau $[u,\ \text{morceau},\ m_1]$ trié par $u$ ; chaque clé est cherchée par dichotomie,
puis \emph{toutes} ses entrées sont lues. Le coût d'une requête est
$O(\nq \log \Nu + A)$, où seul le nombre $A$ d'appariements lus dépend de la taille de la
base. La taille attendue, $\Nu = N_b\,L\,\rho\,F \approx \num{7.5e5}$ entrées pour 100
extraits de \qty{30}{\second}, est proche des \num{677171} mesurées (environ 225 par seconde
de musique, les ancres de fin d'extrait ayant moins de $F$ cibles), soit
\qty{14.6}{\mega\byte}. La \cref{fig-passage-echelle} confirme les deux prédictions : la taille
croît linéairement (de \qty{1.8}{\mega\byte} pour 10 morceaux à \qty{56}{\mega\byte} pour 500)
et le temps de requête médian reste presque constant (\qtyrange{1.9}{2.4}{\milli\second}). Ces temps comptent l'identification à partir du signal prétraité, pas le décodage du
fichier.

\begin{figure}[t]
  \centering
  \inclurefig[0.7\linewidth]{temps_et_taille_index}{temps de requête et taille de l'index en fonction du nombre de morceaux}
  \caption{Passage à l'échelle, de 10 à 500 morceaux indexés. (a)~Temps médian d'une
  requête de \qty{5}{\second} non dégradée : quasi constant. (b)~Taille de l'index en
  mémoire : proportionnelle au nombre de morceaux (tirets : \qty{0.11}{Mo} par morceau).}
  \label{fig-passage-echelle}
\end{figure}

\subsection{Alignement temporel et décision (Q5)}
\label{sec-decision}

\paragraph{Alignement.} Si la requête commence à la trame $m_0$ de la référence, la
propriété (1) assure qu'une paire trouvée dans la requête à l'ancre $m_q$ l'est dans la
référence en $m_r = m_q + m_0$. Tous les appariements corrects vérifient donc
\begin{equation}
  \delta = m_r - m_q = m_0 ,
  \tag{7}
\end{equation}
alors que les appariements fortuits donnent des $\delta$ dispersés. Le score $s$ d'un
morceau est le maximum de son histogramme des $\delta$ ; on répond « le morceau de score
maximal » si $s \ge \sthr$, « inconnu » sinon. Sur la chaîne « jouet »
(\cref{fig-histo-delta}), 699 appariements du bon morceau tombent dans une seule case, en
$\delta$ = \qty{6.66}{\second} pour une position réelle de \qty{6.65}{\second}, alors qu'un
mauvais morceau n'en réunit jamais plus de 2. Les pics secondaires du bon morceau viennent
vraisemblablement d'une boucle musicale répétée.

\begin{figure}[t]
  \centering
  \inclurefig[0.75\linewidth]{histogramme_delta}{histogramme des delta : mauvais morceau contre bon morceau}
  \caption{Histogramme des décalages $\delta = m_r - m_q$ pour une requête de
  \qty{8}{\second} non dégradée (chaîne « jouet »). À gauche, un mauvais morceau : votes
  dispersés, score 2. À droite, le bon morceau : pic net de 699 votes en $\delta = m_0$.}
  \label{fig-histo-delta}
\end{figure}

\paragraph{Score du bon morceau.} Toutes les paires correctes votent dans la même case,
donc d'après~\eqref{eq-cles-correctes}
\begin{equation}
  s_{\text{bon}} \approx F\,\rho\,\Tq\,p^2 ,
  \label{eq-score-bon}
\end{equation}
soit environ 300 pour l'exemple de l'énoncé et $1250\,p^2$ avec nos valeurs
($\rho = 25$, $\Tq = \qty{5}{\second}$). C'est une borne haute : les ancres de fin de requête
n'ont pas toutes leurs cibles ($\nq \approx \num{1015}$ paires et non \num{1250}), et même sans
dégradation $p < 1$, car l'extrait ne tombe pas sur la grille des trames. La
\cref{sec-theorie-mesure} mesure ces deux effets.

\paragraph{Score d'un mauvais morceau.} Si les clés étaient uniformes, les appariements
fortuits avec un mauvais morceau suivraient une loi de Poisson de moyenne
$\nq N_c/M \approx \num{1.1}$ (pour $N_c \approx \num{6800}$ entrées par extrait), répartie
sur environ 550 cases de $\delta$ : un seuil de 3 suffirait à garder la probabilité de fausse
détection sous \qty{1}{\percent} (\cref{eq-plancher}, \cref{sec-annexe-scores}). C'est un
\emph{plancher}. Nous mesurons pourtant des scores négatifs jusqu'à 26 (de 10 à 25 d'après
l'énoncé), car le modèle uniforme suppose deux choses fausses. D'abord, \emph{les clés ne
sont pas uniformes} : la gamme tempérée concentre les pics sur les notes, et les intervalles
consonants reviennent dans tous les morceaux ; il faut remplacer $M$ par un alphabet effectif
$\Meff = 1/\sum_u \pi_u^2$ bien plus petit. Ensuite, \emph{les appariements fortuits ne sont
pas indépendants} : une note tenue produit la même clé à plusieurs ancres consécutives, et une
note tenue commune à la requête et à un mauvais morceau donne plusieurs votes dans la
\emph{même} case, exactement comme un bon morceau. Deux ou trois coïncidences de ce type
suffisent à expliquer des scores de 10 à 25. Le seuil doit donc être calibré sur des
requêtes négatives.

\paragraph{Question 5 : fixer $\sthr$ automatiquement.} On cherche le plus petit seuil qui
garantit un \TFD{} visé $\alpha$ (critère de Neyman-Pearson). La méthode :
(1)~séparer les négatives en un jeu de calibration et un jeu de test, extraits de morceaux
différents ; (2)~ajuster sur les scores de calibration une loi des valeurs extrêmes, la loi
de Gumbel~\cite{gumbel1958}, par la méthode des moments ; (3)~l'inverser pour obtenir le
seuil qui donne $\Prob(S > \sthr) = \alpha$, ce qui extrapole au-delà du plus grand score
observé (\cref{sec-annexe-scores}) ; (4)~valider sur les négatives de test. Cette validation
est limitée par le nombre $n$ de négatives : si aucune n'est acceptée, la borne à
\qty{95}{\percent} du \TFD{} vaut environ $3/n$, soit \qty{6}{\percent} pour $n = 50$ ; il
faudrait 300 négatives sans erreur pour garantir moins de \qty{1}{\percent}. Le modèle indique
aussi que le seuil ne croît qu'avec le logarithme du nombre de morceaux.

\paragraph{Mise en œuvre et résultat.} Le plus fort score des 50 négatives de calibration,
toutes conditions confondues, vaut 23 ; nous avons retenu
$\sthr = \lceil \num{1.5} \times 23 \rceil = 35$. La marge est nécessaire : sur les négatives de
\emph{test}, le score monte jusqu'à 26, ce que seul le découpage calibration/test révèle. Avec
$\sthr = 35$, $\TFD = 0/50$ dans toutes les conditions. La loi de Gumbel, testée, donne
$\sthr \approx 16$ à 18 pour $\alpha = \qty{1}{\percent}$, mais un \TFD{} réel de 4 à
\qty{5}{\percent} (\cref{tab-gumbel}) : sa queue est trop légère, pour les deux raisons
ci-dessus. Elle donne l'ordre de grandeur ; le seuil prudent de 35 garantit l'absence de
fausse détection.

\paragraph{Séparation des scores.} La \cref{fig-histo-scores} montre que les deux
populations sont séparées d'un ordre de grandeur : à $\Tq = \qty{5}{\second}$, les négatives
ne dépassent pas 23, les positives vont pour la plupart de 30 à 900. Les queues se
recouvrent entre 10 et 25 : quelques positives très dégradées (bruit à \qty{0}{\decibel},
réverbération à \qty{1}{\second}) deviennent des non-détections, jamais des confusions. Sur
la \cref{fig-tbd-tfd}, ces deux conditions sont les seules où il faut arbitrer : un seuil
choisi pour elles seules (18 à 20) donnerait 83 et \qty{82}{\percent} de \TBD{} à $\TFD = 0$,
contre \num{66.5} et \qty{51}{\percent} avec le seuil unique, qui doit tenir pour toutes les
conditions à la fois.

\begin{figure}[t]
  \centering
  \inclurefig[0.78\linewidth]{histogrammes_scores}{histogrammes des scores positifs et négatifs, abscisse log, seuil}
  \caption{Histogrammes des scores des requêtes positives et négatives, avec le seuil
  $\sthr = 35$. $\Tq = \qty{5}{\second}$, toutes dégradations réunies (brouhaha compris) :
  \num{3000} requêtes positives et \num{1500} négatives. (a)~Vue d'ensemble en abscisse
  logarithmique ; les scores étant entiers, les classes sont bornées sur des demi-entiers et
  ramenées à une même largeur en log. (b)~Zoom linéaire de 0 à 60 : aucune négative entre
  23 et 35, c'est la marge du seuil.}
  \label{fig-histo-scores}
\end{figure}

\begin{figure}[t]
  \centering
  \inclurefig[0.72\linewidth]{courbe_tbd_tfd}{TBD en fonction du TFD en balayant le seuil, une courbe par dégradation}
  \caption{\TBD{} en fonction du \TFD{}, obtenu en balayant le seuil, pour six conditions,
  $\Tq = \qty{5}{\second}$, $N^+ = 200$. La courbe ne dépend pas du seuil retenu : elle
  utilise les 100 négatives (calibration et test), d'où un \TFD{} qui avance par pas de
  \qty{1}{\percent}. (a) Vue d'ensemble ; (b) zoom sur le coin supérieur gauche, où les courbes
  des conditions faciles se confondent avec le bord. Les cercles marquent le seuil retenu,
  $\sthr = 35$.}
  \label{fig-tbd-tfd}
\end{figure}

\subsection{Récapitulatif des paramètres}
\label{sec-parametres}

Le \cref{tab-parametres} rassemble la valeur retenue pour chaque paramètre et la mesure ou
la propriété qui la justifie.

\begin{table}[!htb]
  \centering\footnotesize
  \caption{Paramètres utilisés pour toutes les mesures, et leur justification. Toutes les
  valeurs sont celles de départ du sujet ; les balayages confirment qu'aucune autre valeur
  testée ne fait mieux.}
  \label{tab-parametres}
  \begin{tabularx}{\linewidth}{@{}llX@{}}
    \toprule
    Paramètre & Valeur & Justification \\
    \midrule
    $\Fe$ & \qty{8}{\kilo\hertz} & Shannon, bande utile \qtyrange{0}{4}{\kilo\hertz} (\cref{sec-shannon}) \\
    Fenêtre & Hann périodique & lobes secondaires faibles ; continue, donc (3) vaut avec une erreur d'ordre $L_h/N$ (\cref{sec-proprietes}) \\
    $N$ & 1024 & sépare les demi-tons dès \qty{190}{\hertz}, $\Tw = \qty{128}{\milli\second}$ (\cref{sec-tfct}) \\
    $H$, $\Nfft$ & $N/2$, $N$ & somme des fenêtres constante ; pas d'interpolation inutile (\cref{sec-tfct}) \\
    $\mv$, $\kv$ & 3, 5 & condition (5) ; avec le quota, $\rho p^2 = \num{7.8}$, le double d'un voisinage sans quota de même $\rho$ (\cref{fig-p-rho}) \\
    Seuil, quota & médiane + \qty{6}{\decibel}, 5 pics/bande/s & invariance au gain, maîtrise de $\rho$ ; $\rho$ mesuré $\approx 25$~pics/s (\cref{sec-pics}) \\
    Bandes & 6, géométriques, \qtyrange{200}{4000}{\hertz} & même intervalle musical par bande (\cref{sec-pics}) \\
    Zone cible & $\dm \in [2, 31]$, $|\dk| \le 60$ & exclut les paires instables ; $\Nt = 30$ (\cref{sec-empreintes}) \\
    $F$ & 10 & $Fp > 1$ ; plus petit $F$ qui atteint le plateau de \TBD{} (5, 10, 15 testés, \cref{sec-empreintes}) \\
    $\gf$ & 1 & même \TBD{} que $\gf = 2$, moins de parasites ; $\gf = 8$ fait chuter le \TBD{} (\cref{tab-gammaf}) \\
    $\Tq$ & \qty{5}{\second} & score $\propto \Tq$ (\cref{eq-score-bon}) ; 3 et \qty{10}{\second} aussi mesurés \\
    $\sthr$ & 35 & $\lceil \num{1.5} \times{}$ plus fort score négatif de calibration$\rceil$ ; $\TFD = 0/50$ en test (\cref{sec-decision}) \\
    \bottomrule
  \end{tabularx}
\end{table}
